\lecture{9}{24 March.}{}
\section{First attempt: Outer Measure}
\begin{definition}
    An open box \(B\) in \(\mathbb{R} ^n\) is any set of the form 
    \[
        B = \prod _{i=1}^n (a_i, b_i) \coloneqq \left\{ (x_1, \dots , x_n) \in \mathbb{R} ^n : x_i \in (a_i, b_i) \text{ for all } 1 \le i \le n  \right\},
    \]  
    where \(b_i \ge a_i\) are real numbers. We define the volume \(\mathrm{vol} (B) \) of this box to be the number 
    \[
        \mathrm{vol} (B) \coloneqq \prod _{i=1}^n (b_i - a_i). 
    \]   
\end{definition}

\begin{definition}
    Let \(\Omega \subseteq \mathbb{R} ^n\) be a subset of \(\mathbb{R} ^n\). We say a collection \((B_j)_{j \in J}\) of boxes cover \(\Omega \) iff \(\Omega \subseteq \bigcup_{j \in J} B_j \).    
\end{definition}

\begin{remark}
    An index set \(J\) is countable if it is finite or \(\exists \) a one-to-one map \(f: J \to \mathbb{N} \) where \(\mathbb{N} = \left\{ 1, 2, \dots  \right\} \) is the set of natural numbers.
\end{remark}

\begin{remark}
    The set of rational numbers is countable. We usually denote it by
    \[
        \mathbb{Q} = \left\{ \frac{p}{q} \mid q \neq 0, p, q \in \mathbb{Z}  \right\}. 
    \]
\end{remark}

\begin{definition}[Outer measure]
    If \(\Omega \) is a set in \(\mathbb{R} ^n\), we define the outer measure \(m^* (\Omega )\) of \(\Omega \) to be the quantity 
    \[
        m^*(\Omega ) \coloneqq \inf \left\{ \sum_{j \in J} \mathrm{vol}(B_j): (B_j)_{j \in J} \text{ is a collection of boxes covers } \Omega \text{ and } J \text{ at most countable.}      \right\}. 
    \]    
    Note that \(m^*(\Omega )\) may be \(\infty \). Also, since \(\mathrm{vol} (B_j) \ge 0 \), so \(m^*(\Omega ) \ge 0\). 
\end{definition}

\begin{lemma}
    The space \(\mathbb{R} ^n\) can be covered by countably many translates of the open cube \(\left( -\frac{1}{2}, \frac{1}{2} \right)^n \). 
\end{lemma}

\begin{proof}
    For each \(m \in \mathbb{Z} ^n\), define 
    \[
        Q_m \coloneqq \frac{m}{2} + \left( -\frac{1}{2}, \frac{1}{2} \right)^n. 
    \] 
    We claim \(\mathbb{R} ^n = \bigcup_{m \in \mathbb{Z} ^n} Q_m \). Let \(x = (x_1, x_2, \dots , x_n) \in \mathbb{R} ^n\). For each \(i\), set \(m_i \coloneqq \lfloor 2x_i \rfloor \in \mathbb{Z} \), then 
    \[
        \frac{m_i}{2} \le x_i < \frac{1}{2} + \frac{m_i}{2} \implies x \in \frac{m}{2} + \left( -\frac{1}{2}, \frac{1}{2} \right)^n. 
    \]
    Note that \(\mathbb{Z} ^n\) is countable since \(\mathbb{Z} \) is countable, and thus we're done.      
\end{proof}

\begin{lemma}[Properties of outer measure]
    Outer measure has the following \(6\) properties:
    \begin{itemize}
        \item [(v)] \textbf{(Empty set)} The empty set \(\varnothing \) has outer measure \(0\), i.e. \(m^*(\varnothing ) = 0\).  
        \item [(vi)] \textbf{(Positivity)} We have \(0 \le m^*(\Omega ) \le \infty \) for all measurable set \(\Omega \).  
        \item [(vii)] \textbf{(Monotocity)}  If \(A \subseteq B \subseteq \mathbb{R} ^n\), then \(m^*(A) \le m^*(B)\).  
        \item [(viii)] \textbf{(Finite sub-additivity)} If \((A_j)_{j \in J}\) are finite collection of subsets of \(\mathbb{R} ^n\), then 
        \[
            m^* \left( \bigcup_{j \in J} A_j  \right) \le \sum_{j \in J} m^*(A_j).  
        \]
        \item [(x)] \textbf{(Countable sub-additivity)} If \((A_j)_{j \in J}\) are countable collection of subsets of \(\mathbb{R} ^n\), then 
        \[
            m^* \left( \bigcup_{j \in J} A_j  \right) \le \sum_{j \in J} m^*(A_j).  
        \] 
        \item [(xiii)] \textbf{(Translation invariance)} If \(\Omega \) is a subset of \(\mathbb{R} ^n\), and \(x \in \mathbb{R} ^n\), then \(m^* (x + \Omega ) = m^*(\Omega )\).    
    \end{itemize} 
\end{lemma}

\begin{proof}[proof of (v)]
    From the definition we know \(m^* (\Omega ) \ge 0\) for all sets \(\Omega \subseteq \mathbb{R} ^n \). Hence, \(m^*(\varnothing ) \ge 0\). Now given \(\varepsilon > 0\), we know \(\left( 0, \varepsilon^{\frac{1}{n}} \right)^n \) is an open box that covers \(\varnothing\). Thus,
    \[
        m^* (\varnothing ) \le \mathrm{vol}\left( \left( 0, \varepsilon ^{\frac{1}{n}} \right)^n  \right) = \varepsilon.  
    \]
    Thus, \(m^* (\Omega ) = 0\).        
\end{proof}

\begin{proof}[proof of (vi)]
    It follows from the definition of \(m^*\). 
\end{proof}

\begin{proof}[proof of (vii)]
    If \(m^* (B) = \infty \), then this is trivial. Now suppose \(m^* (B) < \infty \). Recall 
    \[
        m^* (B) = \inf \left\{ \sum_{j \in J} \mathrm{vol} (B_j) : (B_j)_{j \in J} \text{ are open boxes cover } B    \right\}. 
    \]  
    Also, for any countable collection of boxes covering \(B\), say \((B_j)_{j \in J}\), it also covers \(A\), so
    \[
       \left\{ \sum_{j \in J} \mathrm{vol} (B_j) : (B_j)_{j \in J} \text{ are open boxes cover } B    \right\} \subseteq \left\{ \sum_{j \in J} \mathrm{vol} (A_j) : (A_j)_{j \in J} \text{ are open boxes cover } A    \right\},
    \] 
    which shows
    \[
        m^* (A) \le m^*(B).
    \]  
\end{proof}

\begin{proof}[proof of (x)]
    If one of \(m^* (A_j) = \infty \), then 
    \[
        \infty = m^* (A_j) \le m^* \left( \bigcup_{j \in J} A_j  \right), 
    \] 
    so we're done. Now we assume \(m^* (A_j) < \infty \) for all \(j \in J\). Given \(\varepsilon > 0\), \(\exists \) open boxes \(B_{j, k}\) s.t. \(A_j \subseteq \bigcup_{k=1}^{\infty} B_{j, k} \) and 
    \[
        \sum_{k=1}^{\infty} \mathrm{vol} (B_{j, k}) \le m^* (A_j) + \frac{\varepsilon}{2^j}
    \]
    since the definition of \(m^*\) is infimum. Now since \(J\) is countable, so we can re-label it to \(\mathbb{N}\). Hence, 
    \[
        \bigcup_{j=1}^{\infty} A_j \subseteq \bigcup_{j=1}^{\infty} \bigcup_{k=1}^{\infty} B_{j, k},   
    \]       
    and note that R.H.S. is the union of a collection of boxes, so \(\bigcup_{j=1}^{\infty} A_j \) is covered by them.  
    This shows 
    \begin{align*}
        m^* \left( \bigcup_{j=1}^{\infty} A_j  \right) \le \sum_{j=1}^{\infty} \sum_{k=1}^{\infty} \mathrm{vol} (B_{j, k}) \le \sum_{j=1}^{\infty} m^* (A_j) + \frac{\varepsilon}{2^j} = \sum_{j=1}^{\infty} m^* (A_j) + \varepsilon.     
    \end{align*}
    Now since this is true for all \(\varepsilon > 0\), so 
    \[
        m^* \left( \bigcup_{j=1}^{\infty} A_j  \right) \le \sum_{j=1}^{\infty} m^*(A_j).  
    \] 
\end{proof}

\begin{proposition}[Outer measure of closed box]
    For any closed box
    \[
        B = \prod _{i=1}^n [a_i, b_i] = \left\{ (x_1, x_2, \dots , x_n) \mid a_i \le x_i \le b_i \text{ for } 1 \le i \le n  \right\},
    \]
    we have
    \[
        m^* (B) = \prod _{i=1}^n (b_i - a_i).
    \]
\end{proposition}

\begin{proof}
    Write 
    \[
        \mathrm{vol} (B) = \prod _{i=1}^n (b_i - a_i). 
    \]
    We first show that 
    \[
        m^* (B) \le \mathrm{vol}(B). 
    \]
    Let \(\varepsilon > 0\), then the box 
    \[
        U_{\varepsilon } \coloneqq \prod _{i=1}^n (a_i - \varepsilon, b_i + \varepsilon) 
    \] 
    contains \(B\). Thus,
    \[
        m^* (B) \le \mathrm{vol}(U_{\varepsilon} ) = \prod _{i=1}^n (b_i - a_i + 2\varepsilon ). 
    \] 
    Take \(\varepsilon \to 0\), then we can show 
    \[
        m^* (B) \le \mathrm{vol} (B). 
    \] 
    Now we show 
    \[
        m^* (B) \ge \prod _{i=1}^n (b_i - a_i).
    \]
    \begin{claim}
        For any countable family of open boxes \(\left\{ B_j \right\}_{j \in J} \) covering \(B\). We have
        \[
            \mathrm{vol} (B) \le \sum_{j \in J} \mathrm{vol}(B_j).
        \]   
    \end{claim}

    \begin{explanation}
        Since \(B\) is compact, \(\exists \) a finite subcover of \((B_j)_{j \in J}\)
        \[
            \left( U_j \right)_{j=1}^N \text{ s.t. } B \subseteq \bigcup_{j=1}^{N} U_j.    
        \]
        Thus,
        \[
            \sum_{j=1}^N \mathrm{vol} (U_j) \le \sum_{j \in J} \mathrm{vol} (B_j).    
        \]
        Choose a closed box \(Q\) s.t. \(B \subseteq Q\) and \(U_k \subseteq Q\) for \(1 \le k \le N\). Define functions on \(Q\) by 
        \[
            f(x) = \begin{dcases}
                1, &\text{ if } x \in B ;\\
                0, &\text{ if } x \in Q \setminus B.
            \end{dcases} \quad f_k(x) = \begin{dcases}
                1, &\text{ if } x \in U_k ;\\
                0, &\text{ if } x \in Q \setminus U_k .
            \end{dcases}
        \]  
        Thus,
        \[
            \int _Q f(x) \,\mathrm{d} x = \mathrm{vol} (B) \text{ and } \int _Q f_k(x) \, \mathrm{d} x = \mathrm{vol} (U_k).     
        \]
        Note that this is integrable since \(B, \left( U_k \right)_{k=1}^N \) are all boxes and \(Q\) is also a box. 
        Now since 
        \[
            f(x) \le \sum_{j=1}^k f_j(x) \text{ on } Q.
        \] 
        (If \(x \in B\), then \(x \in U_j\) for some \(j\).)   
        Thus,
        \[
            \mathrm{vol}(B) = \int _Q f(x) \, \mathrm{d} x \le \int_Q \sum_{j=1}^k f_j (x) \, \mathrm{d} x = \sum_{j=1}^k \int_Q f_j (x) \, \mathrm{d} x = \sum_{j=1}^k \mathrm{vol} (U_j) \le \sum_{j \in J} \mathrm{vol} (B_j) .     
        \]
    \end{explanation}

    This implies \(\mathrm{vol} (B) \) is a lower bound of 
    \[
        \left\{ \sum_{j \in J} \mathrm{vol}(B_j) : (B_j)_{j \in J} \text{ are boxes covering } B    \right\}, 
    \] 
    so by the definition of \(m^*\),
    \[
        m^* (B) \ge \mathrm{vol} (B). 
    \] 
\end{proof}